%Dette er en helt simpel skabelon, som kan bruges til at skrive større opgaver og rapporter i. Koden er bevidst lavet så ukompliceret som muligt,
%så dokumentet samtidigt kan bruges som en indgang til at læse LaTeX at kende.

%---------------------------------------------------------
%Dette er kun preample. Det behøver du ikke bekymre dig om. Går ned til der står "Start her". Eller til der hvor sidehovedet indstilles.
%---------------------------------------------------------
\documentclass[a4paper,12pt,oneside]{memoir}
% lidt marginer
\setlrmarginsandblock{3cm}{*}{1}
\setulmarginsandblock{3cm}{*}{1}
\checkandfixthelayout[nearest]
\raggedbottom
\setlength{\parindent}{0pt}%Fjerner indrykning ved nyt afsnit

% Standard dansk opsætning
\usepackage[utf8]{inputenc}   %æøå med utf8-encoding
\usepackage[danish]{babel}            % dansk opsætning
\renewcommand\danishhyphenmins{22}    % fikser en fejl i babel
\usepackage[T1]{fontenc} %fontencoding
\usepackage{lmodern} %sætter skrifttypen til Latin Modern

% Matematiske symboler og fede tegn i ligninger
\usepackage{amsmath, amssymb, bm, mathtools, mathdots}

%Flere mulighder for understregning
\usepackage[normalem]{ulem}

% Tabeller og søjler
\usepackage{array, booktabs, tabularx}

% Figurer og farver
\usepackage{graphicx, caption, subfig, xcolor}
\captionsetup{font=small,labelfont=bf}

%--------------------------------------------------------
%Her sættes sidehoved og sidefod. Skriv dit eget navn ind.
%--------------------------------------------------------
\makepagestyle{mypagestyle} %Min pagestyle med sidetal
\copypagestyle{mypagestyle}{empty}
\makeoddhead{mypagestyle}{Navn\\Klasse}{\quad}{\today\\Fag}
\makeheadrule{mypagestyle}{\textwidth}{\normalrulethickness}
\makeoddfoot{mypagestyle}{}{\thepage}{} %Kræver to oversættelser
\aliaspagestyle{chapter}{mypagestyle}
\pagestyle{mypagestyle} %aktiver ny sidehoved/-fod

\chapterstyle{section}

%Punktopstilling
\usepackage{enumerate}

%Enheder og tal
\usepackage[output-decimal-marker={,}]{siunitx}
\usepackage{icomma} %Komma som decimalseperator

%Talmængder
\newcommand{\R}{\mathbb{R}} %\R for reelle tal
\newcommand{\C}{\mathbb{C}} %\C for komplekse tal
\newcommand{\N}{\mathbb{N}} %\N for naturlige tal
\newcommand{\Z}{\mathbb{Z}} %\Z for hele tal
\newcommand{\Q}{\mathbb{Q}} %\Q for rationale tal
\newcommand{\dg}{^{\circ}}%Brug \dg til at indsætte gradtegn

%Til hyperlinks
\usepackage[hidelinks=true]{hyperref}
\usepackage{memhfixc}

\begin{document}%Her begynder indholdet af dokumentet.
\frontmatter
%-------------------------------------------
%START HER!
%-------------------------------------------

%Forside
\begin{titlingpage}
\thispagestyle{empty}
\vspace*{3cm}
\centering
\Huge \bfseries \sffamily Skabelon til større projekter\\

\vspace*{3cm}
\includegraphics[width=7cm]{cirkler}

\vspace*{4cm}
\begin{flushright}
\large
Lavet af: Navn\\
Klasse: 2.x\\
Fag: Matematik\\
Dato: \today
\end{flushright}
\end{titlingpage}

%Indholdfortegnelse
\tableofcontents*

\mainmatter
\chapter{Første kapitel}
Her er teksten til det første kapitel. Den kunne indeholde mange interessante ting.
\section{Afsnit}
Der er sørme også et afsnit!

\chapter{Andet kapitel}\label{kapitel2}
Her er et andet kapitel, som altid starter på en ny side. Her er det side \pageref{kapitel2}.

\chapter{Tabeller og figurer}

\section{Tabeller}
Hvis man vil lave en simpel tabel, kan man gøre det på følgende måde
\begin{table}[htb]
\centering
\begin{tabular}{|l|r|}
\hline 
$n$ & Primtal mindre end eller lig $n$ \\ 
\hline 
1 & 0 \\ 
\hline 
2 & 1 \\ 
\hline 
3 & 2 \\ 
\hline 
4 & 2 \\ 
\hline 
\end{tabular}
\caption{Primtal mindre end eller lig et naturligt tal.}\label{tabel-primtal} 
\end{table}

Den tabel har nummeret \ref{tabel-primtal}. Tabellen er indlagt som en såkaldt \emph{float}, så \LaTeX\ sætter den selv ind på det sted i teksten, hvor den synes, der er bedst plads. Så behøver man ikke bekymre sig om det.

\section{Figurer}
Man kan også bruge floats til figurer. Det har vi gjort for figur \ref{fig-cirkler}.

\begin{figure}[htb]
\centering
\includegraphics[width=5cm]{cirkler}
\caption{To cirkler.}\label{fig-cirkler}
\end{figure}

Man kan læse mere i \cite[kapitel~6]{latex-bog}, hvor brugen af floats er beskrevet i langt større detalje.

\backmatter
%Litteraturliste
\begin{thebibliography}{1}
\bibitem{latex-bog}
Lars Madsen.
\textit{Introduktion til \LaTeX{}}.
2011.
\bibitem{anden-bog}
Jens Carstensen.
\textit{Komplekse tal}.
Systime,
3. udgave, 2. oplag, 1996.
\end{thebibliography}
%---------------------------------------------------------
%Dokumentet er slut.
%---------------------------------------------------------
\end{document}%Alt efter denne linje kommer ikke med i dokumentet.